\documentclass[11pt,a4paper]{article}
\usepackage[T1]{fontenc}
\usepackage{lmodern,amsmath,amssymb,booktabs,microtype}
\usepackage[margin=24mm]{geometry}
\usepackage[colorlinks=true,urlcolor=blue,linkcolor=blue]{hyperref}
\newcommand{\EE}{\mathbb E}
\newcommand{\tr}{\operatorname{tr}}
\newcommand{\W}[1]{\mathsf W_{#1}}
\newcommand{\Db}{\overline D}
\newcommand{\J}{J_s}
\title{Mixed spectra of the first memory fields\\An upper-charge lift and a positive three-field measure}
\author{Alan Day \and Chengheng\thanks{AI research assistant (OpenAI).}}
\date{Local research note --- 6 November 2026}
\begin{document}\maketitle
\begin{abstract}
We close the two-energy cross measure between the first memory
tangent and the projected second memory coefficient. The cubic
residual has vanishing cross with the creation-force direction,
as proved by a complete endpoint certificate. A direct commutator
lift handles the upper-charge sector in its cross with the linear
rank-filtered field; the same signed-covariance argument extends
to the quartic residual. These results yield a positive spectral
matrix for three memory fields and the covariance of an explicitly
projected quadratic polynomial in memory time.
\end{abstract}
\section{The completed mixed object}
We use standard CAR, the full overlapping complex Gaussian
$p$-body ensemble, $h=H_m$, fixed half filling $m=N/2$, normalized
maximally mixed trace, and annealed inner product
$\langle U,V\rangle=\EE\tr_m U^\dagger V$. Set $\hbar=1$ and take
\begin{equation}
 \lambda_N=p^2/N\longrightarrow\lambda>0,\qquad
 q=e^{-4\lambda},\qquad C_N=\EE\tr_mh^2\longrightarrow C=\J^2/\lambda.
\end{equation}
The fields $Y=Y_-$, $R$, $F$, and $Q_4$ are those of the preceding
notes~\cite{local}. In particular, with $J_hU=\{h,U\}/2$,
\begin{equation}
 L_N=R_N+aJ_hY_N+F_N+o_{L^2}(1),\qquad
 a=-\frac{(1-q)^2}{1+q^2}.
\end{equation}
$L_N=mZ_1-C_{1,m}$ is the conserved-subtracted first memory
coefficient. Here $L$ names a field, not the commutator with $h$.
For the projected second coefficient,
\begin{equation}
 X_N=(1-\widehat\Pi_N)mZ_2
     =Q_{4,N}+8\lambda_N C_Nq_N^2T_N+o_{L^2}(1),
\end{equation}
where $\widehat\Pi_N$ projects onto
$\operatorname{span}\{I,B,O_-,J_hF,E_0,M_0,S_0,Z_0\}$.
The main result is the full two-energy statement
\begin{equation}
 \boxed{\ d\nu_{XL}=d\nu_{XY}=0.\ } \label{main}
\end{equation}
The $Y$ cross was established previously; this note supplies the
three missing ingredients for the $L$ cross. Multiplication by
any fixed polynomial of the left and right energies preserves
these zero cross measures.

\newpage
\section{Creation force versus the cubic residual}
Write $d_i=[H,c_i]$, $A_- =\sum_i d_i^\dagger d_i$,
$A_+=\sum_i d_id_i^\dagger$,
$O_\pm=\W2 A_\pm$, $B=\W2h^2$, and
\begin{equation}
 Y_\pm=O_\pm-cB,\qquad c=\frac{2\lambda(1-q)}{1+q}.
\end{equation}
All $\W j$ in this note are coupling-coordinate Gaussian chaos.
The earlier strong output-deficit bridge gives
\begin{equation}
 T_N=\frac{Y_-+Y_+}{2}-\rho Y_-+o_{L^2}(1),\qquad
 \rho=\frac{v}{v+d_v},
\end{equation}
with $v=2\lambda(1-q)+16\lambda^2q/(1+q)$ and $d_v=2\lambda(1+q)$.
Thus $\nu_{TR}$ reduces to the two oriented force crosses.

For the cubic field put $D_3=(1+q^2)(1+q+q^2)$. In the
lexicographic three-slot mark order $01,02,12$, the numerator
weights of $R$ are
\begin{equation}
 [-q(1+q+q^2),\ D_3,\ -q(1+q+q^2)],
\end{equation}
and its unmarked numerator is $-2\lambda(1-q)(1+q^3)$.
For $Y_+$ the two-slot mark numerator is $1+q$ and the unmarked
numerator is $-2\lambda(1-q)$, with denominator $1+q$.

Take a cubic normal mark $(a,b)$ and a quadratic reverse mark
$(c,d)$. For Gaussian pairing $\pi$, the established local rule is
\begin{equation}
 \alpha=2-4\mathbf1_{\rm incident\ crossing},\qquad
 \delta_{-+}=-4\mathbf1_{\pi(a)=c,\,\pi(b)=d}.
\end{equation}
The quadratic coefficient is $\alpha\beta$ in both orientations;
the equal-orientation linear coefficient is instead
$\delta_{--}=4\mathbf1_{\pi(a)=d,\pi(b)=c}$.
After the Gaussian chaos restrictions, the three- and two-slot
blocks share $k=0,1,2$ bridges. Classify the remaining endpoints
by their $A/B$ colors and the split between exterior arcs. The
number of templates is
\begin{equation}
 \sum_{k=0}^2\binom{5-2k}{3-k}(6-2k)=74.
\end{equation}
Each contains $3!2!/k!$ pairings, for 876 words altogether.
Every coefficient of $\lambda$ and $\lambda^2$ vanishes as an
integer polynomial in $q$, for either force orientation.
The code also checks all 2628 local two-mark minors per orientation
against the determinant implementation.

Exterior spectator chords contribute the same factor on each
endpoint orbit. The certificate therefore covers arbitrary fixed
two-arc energy degrees, not just bounded word lengths. Positive
matrix-measure compactness and moment uniqueness, using the
established diagonal kernels, give
\begin{equation}
 \boxed{\ \nu_{R,Y_+}=\nu_{R,Y_-}=\nu_{R,T}=0.\ } \label{Rcross}
\end{equation}

\newpage
\section{The filtered edge must also reach the upper charge}
Write the neutral rank decomposition $h=\sum_a h_a$, with
$\EE\tr_m h_a^2=C_N\omega_a^{(p)}$, and put $u=p-a$.
The earlier normalized linear residual satisfies
\begin{align}
 f_N&=F_N/\sqrt{V_{F,N}},&
 \widehat f_N&=\sqrt{\lambda_N}\sum_a[R_1(p-a)-1]h_a/\sqrt{C_N},\\
 R_1(u)&=\frac{u(N-2p+u+1)(N-m+p)}{(m-p+1)p^2},&
 \|f_N-\widehat f_N\|_2&\longrightarrow0. \label{fbridge}
\end{align}
This uses the exact rank-weight ratio, factorial input tails, and
the Poisson deficit limit; it is an operator $L^2$ replacement.
A Wick edge from $\widehat f_N$ is
$\sqrt{\lambda_N}(\mathsf K_{p-1}-\mathsf K_p)$ at charge $m$.

Let $U$ be a $(p-1)$-body ensemble normalized at $m$, and set
\begin{equation}
 d=m-p+1,\qquad g(n)=\frac{n-p+1}{d},\qquad \widetilde U=g(Q)U.
\end{equation}
For every allowed charge $n$, the exact reduced rank factors obey
\begin{equation}
 g(n)\frac{\binom{n-a}{p-1-a}}{\binom{m-a}{p-1-a}}
 =\frac{\binom{n-a}{p-a}}{\binom{m-a}{p-a}}.
\end{equation}
Together with $R_1$, this extends the signed covariance to
$m\leftrightarrow n$. The new finite tensor check explicitly
includes $n=m+1$, which the creation-force vertex requires.

Define bilinear commutator marks for neutral $U,V$:
\begin{equation}
 \mathcal M_-(U,V)=-\sum_i[U,c_i^\dagger][V,c_i],\qquad
 \mathcal M_+(U,V)=-\sum_i[U,c_i][V,c_i^\dagger].
\end{equation}
For $D V=QV-\sum c_i^\dagger Vc_i$ and
$\Db V=(N-Q)V-\sum c_iVc_i^\dagger$, direct CAR gives at charge $m$
\begin{align}
 \mathcal M_+(gU,V)&=(1+d^{-1})\mathcal M_+(U,V)-d^{-1}U\Db V,\\
 \mathcal M_+(V,gU)&=(1+d^{-1})\mathcal M_+(V,U)-d^{-1}(\Db V)U,\label{rev}\\
 \mathcal M_-(gU,V)&=(1-d^{-1})\mathcal M_-(U,V)+d^{-1}UDV,\\
 \mathcal M_-(V,gU)&=(1-d^{-1})\mathcal M_-(V,U)+d^{-1}(DV)U.
\end{align}
These formulas preserve operator ordering. For example
$[gU,c]=g[U,c]-d^{-1}cU$; the sum in the remaining term reduces
to $(\Db U)V-\Db(UV)$, proving the first equation.
For a $p$-body $V$, $DV=pV$ and $\Db V=pV-\Delta_V$.
At half filling $\Delta_V$ has bounded fixed moments by the
previous input-deficit formula. Hence every correction in
\eqref{rev} is $O_{L^2}(p/N)=O_{L^2}(p^{-1})$, also inside each
fixed-degree trace by higher-moment bounds.

\newpage
\section{Why the filtered crosses vanish at all energy orders}
It is essential to apply the preceding identities to the actual
reverse commutator mark. The general superoperator formula is
\begin{equation}
 \Db[f(Q)B]_m=f(m+1)(\Db B)_m
              +(N-m)[f(m)-f(m+1)]B_m.
\end{equation}
Replacing a normal body-count mark mechanically by $\Db$ would
leave a correction of order one. Equation~\eqref{rev} includes
the needed terms directly and has the smaller error just proved.

For normal marks in $Q_4$, the old lifting identity extends to
every fixed product degree $k$. If $B$ contains one $(p-1)$-body
factor and $k-1$ factors of body order $p$, then
\begin{equation}
 [(kp-D)(gB)]_m
 =(1-d^{-1})[(kp-1-D)B]_m+\frac{(k-1)p}{d}B_m. \label{quartic}
\end{equation}
In particular $k=4$ treats the quartic vertex. Nonadjacent pair
marks can be written as finite sums of total marks on contiguous
subproducts, so~\eqref{quartic} also treats each of its six pair
positions. If the shifted slot lies outside the mark, the scalar
lift evaluates to one at the exterior charge. No growing number
of factors is used.

The mixed-body chord limit in Ref.~\cite{BNR} applies to index
sizes $p$ and $p-1$: their pair-overlap means have the same
limit, and triple overlaps vanish. The project's marked
determinant calculation, together with absolute exponential
overlap bounds, permits fixed mark derivatives of this limit.
The half-filled trace bridge transfers the argument to the
present state. Consequently each fixed marked graph has the
same leading value when one covariance edge is shifted from
$p$ to the lifted $p-1$ ensemble. Equations~\eqref{rev} and
\eqref{quartic} control the additional vertex terms.

For $B_N=Y_{+,N}$ or $Q_{4,N}$ and fixed $r,s$, replacing $F_N$
by its normalized proxy in~\eqref{fbridge} is legitimate because
\begin{equation}
 |\langle(f_N-\widehat f_N)h^rB_Nh^s\rangle|
 \leq\|f_N-\widehat f_N\|_2
       \big(\EE\tr_m B_Nh^{2r}B_Nh^{2s}\big)^{1/2}\longrightarrow0.
\end{equation}
The second factor is bounded by the already established diagonal
energy moments. Gaussian centering deletes internal pairings of
the $Y_+$ block; $\W4$ deletes those of the quartic block.
Each surviving Wick graph has exactly one signed filtered edge.
Its two lifted terms have identical limits and cancel. Thus all
fixed mixed moments vanish. The positive spectral matrix,
diagonal high moments, and compact candidate support then imply
\begin{equation}
 \boxed{\ \nu_{F,Y_+}=\nu_{F,Q_4}=0,\qquad \nu_{F,T}=0.\ } \label{Fcross}
\end{equation}
The last statement uses the already known $\nu_{F,Y_-}=0$.
Equations~\eqref{Rcross} and~\eqref{Fcross}, the old
$\nu_{Q_4,R}=0$, and the preceding $\nu_{X,Y_-}=0$ prove~\eqref{main}.

\newpage
\section{A positive three-field matrix and its projected polynomial}
Let $\mu_r$ be the existing normalized Mehler two-energy measure,
with the variance-$C$ energy marginal, and put $e=(E+E')/2$.
The complete matrix for $(Y,L,X)$ is
\begin{equation}
 d\boldsymbol\nu=
 \begin{pmatrix}
 W\,d\mu_{q^2}&aeW\,d\mu_{q^2}&0\\
 aeW\,d\mu_{q^2}&d\nu_L&0\\
 0&0&d\nu_X
 \end{pmatrix}, \label{matrix}
\end{equation}
where
\begin{align}
 d\nu_L&=V_Fd\mu_q+V_Rd\mu_{q^3}+a^2e^2W d\mu_{q^2},\\
 d\nu_X&=V_4d\mu_{q^4}+V_2d\mu_{q^2},\\
 W&=4\J^4\left(\lambda^{-1}+\frac{4q}{1+q}\right),\qquad
 V_F=64\J^6q^2,\\
 V_R&=\frac{C^3}{D_3^2}(\lambda A_1+\lambda^2 A_2),\\
 A_1&=4(1-q^2)(1+q^2)(1+q+q^2)^2,\\
 A_2&=16q^2(1+q)(1+q^2)(1+q+q^2)(2q^2-q+2).
\end{align}
The previous note gives the positive $V_4,V_2$ explicitly.
Relative to the common product energy marginal, write
$K_j=d\mu_{q^j}/d(\nu_q\otimes\nu_q)$. Then~\eqref{matrix}
has the exact positive factorization
\begin{equation}
 M\,\operatorname{diag}(WK_2,V_FK_1,V_RK_3,V_4K_4,V_2K_2)M^T,
 \quad
 M=\begin{pmatrix}1&0&0&0&0\\ae&1&1&0&0\\0&0&0&1&1\end{pmatrix}.
\end{equation}
The determinant is positive:
\[
 WK_2(V_FK_1+V_RK_3)(V_4K_4+V_2K_2)>0.
\]
The five fields $Y,T,F,R,Q_4$ have diagonal
two-point measures; $Y$ and $T$ share the normalized $q^2$ kernel.

Define the explicitly filtered polynomial
$\Phi(s)=Y+isL-s^2X/2$. Its energy-resolved covariance is
\begin{align}
 d\nu_{\Phi(s),\Phi(u)}={}&
 W[1+iae(u-s)+a^2e^2su]d\mu_{q^2}\notag\\
 &+su(V_Fd\mu_q+V_Rd\mu_{q^3})
 +\frac{s^2u^2}{4}d\nu_X. \label{poly}
\end{align}
At ordinary center time $T$, reflection $(E,E')\mapsto(-E',-E)$
preserves the frequency and reverses $e$. Thus
\begin{equation}
 \langle\Phi(s)^\dagger\Phi(u;T)\rangle
 =WG_{q^2}(T)+suC_L(T)+\frac{s^2u^2}{4}C_X(T).
\end{equation}
Since $X$ has an additional, partly nonconserved projection,
$\Phi$ is not the raw memory Taylor polynomial. Formula~\eqref{poly}
does not assert a finite-memory remainder bound. Such a bound must
also restore the removed dynamic descendants.

\newpage
\section{Checks, scope, and the next reconstruction step}
\paragraph{Endpoint certificate.}
Both orientations were computed, rather than extrapolating the
new result from the old normal one: 74 templates and 876 words
per orientation, 5256 determinant-minor comparisons in total.
All integer polynomial arrays vanish. The direct certificate
also gives a shorter route to the previously known $R$--$Y_-$
cross; it is used here as a control, not claimed as a new kernel.

\paragraph{Upper-charge and operator checks.}
The exact rank-Casimir covariance tensors were evaluated for
$(N,p,m)=(4,2,2),(5,2,3),(6,2,3)$, now explicitly including
$m+1$. Their 93076 entries agree after clearing denominators.
The previous normal quadratic/cubic controls were retained.
Independent integer Hermitian samples test both locations of a
lifted factor in each CAR orientation, the reverse superoperator
identity, and all four locations in a noncommuting quartic product:
81 exact identities, with 80 nonzero witnesses. Cleared entries
remain below $2^{48}$. No Gaussian sampling is used to estimate
ensemble expectations; the samples only check operator identities.

The first checker required every witness to be nonzero. It stopped
on the auxiliary $\Db(gU)$ identity for one traceless one-body
sample at $(6,2,3)$. Here $\Db U=U-\tr(J)I$ and its cleared
identity reduces to $-3\tr(J)I$, so zero is correct. The original
checker and diagnosis are preserved. The revised checker records
this zero and still requires all reverse/normal commutator and
quartic witnesses to be nonzero. No algebraic tolerance was relaxed.

\paragraph{Assembly checks.}
An independent symbolic multiplication verifies the positive
factorization, determinant, all nine polynomial coefficients in
$(s,u)$, the conjugation symmetry, and the energy-reflection
reduction. Existing diagonal-kernel calculations supply the time
kernels; no new diagonal numerical scan is needed. Every new
scientific output records its inputs, hashes, and environment.

\paragraph{What follows.}
The missing mixed spectra of the specified three fields are now
closed at fixed half filling. The next work must extend the
projection and bounds uniformly in a central charge window,
restore the removed dynamic descendants when constructing the
actual second memory coefficient, and control the two-parameter
finite-memory remainder. The resulting kernel, rather than the
projected polynomial alone, is the object to connect to the
Susskind singlet/meson spatial reconstruction program.
Two-point spectral orthogonality does not establish mixed
higher-point independence or a bulk mass dictionary.

\begin{thebibliography}{9}\small
\bibitem{BNR} M. Berkooz, V. Narovlansky and H. Raj,
\emph{Complex Sachdev--Ye--Kitaev model in the double scaling limit},
\href{https://arxiv.org/html/2006.13983v2}{arXiv:2006.13983v2},
Sec.~3.3, Eqs.~(3.11)--(3.13): oriented words with generic index-set lengths.
\bibitem{TT} J. Tall and S. Tomsovic,
\emph{Double-scaled bosonic and fermionic embedded ensembles,
complex SYK, and the dual Hilbert space},
\href{https://arxiv.org/html/2604.14522v2}{arXiv:2604.14522v2}.
Existing spectral Wick and strength-density framework.
\bibitem{local} A. Day and Chengheng, local project notes (2026-11-06):
\emph{Closing the cubic singlet two-point spectrum};
\emph{Energy-resolved memory background and first fluctuations};
\emph{Two orientations of the fermionic force}.
The corresponding source, exact certificates, and PDFs remain in the project.
\end{thebibliography}
\end{document}
